\section{The Word Remains Within}

American Christians generally use ``logos'' (if they use it) in the sense that Aristotle wants to use it:

\begin{quotex}
For Aristotle, \emph{logos} is something more refined than the capacity to make private feelings public: it enables the human being to perform as no other animal can; it makes it possible for him to perceive and make clear to others through reasoned discourse the difference between what is advantageous and what is harmful, between what is just and what is unjust, and between what is good and what is evil. (PA Rahe)
\end{quotex}


That is, when modern Christians use the word Logos, they predominantly conceive of ``that which makes for argument, or logic''. This definition is by and large the regnant one, even among the classical Christian movement. Practically speaking, the Logos is thought of in terms of that which ``makes sense'' out of our logic, and connects us, through dialectic, to the mind of God.

Although I would not want to argue that this is not a dimension of Logos, I do say that it is an impoverished one. Philo, for instance, terms it the \emph{logos endiathetos (the word remaining within),} while the Stoics (who influenced not to much St. John as Justin the Martyr) believed that the Logos was the generative principle of the universe. Philo, also, assigned to the Logos a kind of demi-urgic status. Surely St. Paul must have had something like this in mind when he described Christ as being the vessel and sum of the worlds, which He would recapitulate and return to the Father at the end of time.

The point here is that Logos is not merely an Aristotelian or Enlightenment concept of active, dialectical Reason, a spark of divinity that ensures our divine image and reflects some logical function of Christ's ontological status. This would be to mechanize the Logos, or at least, trend in this direction. It would be to make the Logos merely a pattern filtered through right Reason (which, of course, it partly is).

Justin Martyr wrote the following:

\begin{quotex}
I shall give you another testimony, my friends, from the Scriptures, that God begot before all creatures a Beginning, who was a certain rational power proceeding from Himself, who is called by the Holy Spirit, now the Glory of the Lord, now the Son, again Wisdom, again an Angel, then God, and then Lord and Logos.
\end{quotex}


The Orthodox want to argue that this ``angel of the Lord'' is actually the divine energies which one experiences in illumination (after purification). It is not a created being or merely external theophany (as Augustine is alleged to have thought) which is brought into and out of existence in order to create a saga of revelation, but rather, participation in the divine energies (the divine essence being reserved for God, forever, alone). ``Being in the Spirit'' allows one to experience this divine energy.

If we take our hint from the Orthodox, we might try thinking of the Logos as \emph{all higher states of being} whatsoever (the lower ones being created by natural deprivation or distortion of the collateral states associated with Logos). The Logos is ``Light'' or ``Life'' -- the ``Kingdom of God''. Understood this way, Jesus was Himself the embodiment in full of the reign of God: He did not exhaust, but rather, fully expressed, the Logos.

The Logos is therefore a spiritual state of Being, a higher kingdom of existence, a noetic faculty of soul, the seeds of Creation, the divine image within, Right Reason, the Nous itself and all higher worlds (excluding God's essence), the Logos Spermatikos and First Born before all Worlds, the Alpha and Omega and the end of time, and the way, truth, and life, which is the light of men.

It is clear that the emphasis in the West has been, for centuries now, upon the dialectical apparatus of the logical mind, and while divine certainty has been ascribed strongly to this (how else could David Hume die in such peace of mind?), it has tended to both arrogate and impoverish the depths and even the width of the human spirit and soul. Christianity itself has been implicitly guilty in this, because Western Christianity tends to think of God, when it does so at all, as a logically accessible entity that operates according to the strictures of common, rational thought; in practice, it is often reducible to the ``structures'' of thought.

With this is mind, next week we will begin our trip through Iamblichus' arithmetic theology, showing how a ``mysticism'' of the numbers is possible that is revelatory of fundamental patterns in both lower and higher reality.


\flrightit{Posted on 2012-12-20 by Logres}

\begin{center}* * *\end{center}

\begin{footnotesize}
\begin{sffamily}
\texttt{Dominion on 2012-12-25 at 03:43 said:}

If the Logos is all higher states of being, one must consider what is to be made of the other two Divine persons of the Trinity. Might we think of the Holy Spirit as the Wisdom that the initiate seeks? The Holy Spirit is called ``the Helper''; Beatrice helps Dante in his journey, and the perennial Sophia who is the true love of the initiate is essential for leading him to God, who reveals himself slowly. Thus we have the Logos (the states of Being) on the one hand, and the Holy Spirit (helper and guide) on the other. In the Christian Tradition, both stem from the Father, the person of the Trinity who ``generates'' (but does not create) the other two. It may be possible to consider the Father the face of the supreme Reality, the fundamental essence of which the states of Being and the Divine Wisdom are reflections, extensions and manifestations.

In this line of thought, the reconciliation of Christ's statement that ``none can come to the Father but through me'' and the transcendent unity of religions taught by the perennial Tradition logically follow.

\hfill

\texttt{Logres on 2012-12-25 at 11:14 said:}

Good thoughts, Dominion. I've heard the Son and Spirit called ``the two hands of the Father'', but (of course) this slides towards seeing them as One Person. Nevertheless, perhaps the unity of God needs more emphasis in our day of fragmentation and needless accretion of knee-jerk thought patterns within Christianity. If Christianity, instead of being seen as a competing religion to supersede all others, was actually a purifier and exalter and perfecter of what was good in those religions, this would be the counterpart of emphasizing God's unity. The Church was often more anti-pagan than pro-Christ, and we're still paying for it. Have you seen Bulgakov's work on Sophia?
\url{http://en.wikipedia.org/wiki/Sergei\_Bulgakov}

It sounds as if you had... .

\hfill

\texttt{Dominion on 2012-12-25 at 13:31 said:}

I have not, but thank you for the link. Your comment on the Church's anti-Heathen vs pro-Christ attitude is apt. Unfortunately, the same attitude has been taken between Christians themselves (the miaphysite question between Eastern and Oriental Orthodox and the Papal issue and filioque between East and West come to mind). Personal ego and prestige have often outdone impersonal and rational dialogue and common cause in the search of truth, right down to when St. Nicholas punched out Arius at Nicaea. If Christians cannot interact with each other, then dialogue on Tawhid and the one-ness of the Divine reality may be far off yet.

\hfill

\texttt{Graham on 2012-12-27 at 09:07 said:}

Gentlemen, on the need for pro-Christ `positive formulations' and `dialogue' ... did the last five decades never happen?

``Where there is no hatred of heresy, there is no holiness.''

Dialogue is `vewwy nice,' but has to be sharply qualified in terms of its final cause. On the exoteric level, it must aim at conversion; on the esoteric, spiritual depth; on all levels, at the correction of error. When the final cause is ignored, dialogue and positive formulations become in effect ends in themselves, and thereby tools of the Demon sowing confusion.

\hfill

\texttt{Dominion on 2012-12-27 at 13:30 said:}

That's more or less my point. Much of the dialogue that has taken place has been in order to render reconciliation on the exoteric plane at the expense of Truth. This is especially true with the Catholic Church. That said, esoteric dialogue has also occurred and indeed been fruitful. The Traditionalist writers are fine examples of that.

\hfill

\texttt{Jason-Adam on 2012-12-27 at 17:35 said:}

The problem I have with hesychasm is that it is Orthodox and Orthodox practices can not be adopted by Catholics without sacrificing Catholicity -- despite what one wants to think, there are huge differences between the Western and Eastern Churches that ultimately come down to culture. When all is said and done I think Orthodox Christianity is as foreign to the West as Islam or Hinduism is despite the same Triune Deity we worship. Remember the Tradition is one but the forms are many.

I used to be very interested in Western Rite Orthodoxy, until I found that becoming Orthodox basically involves adopting an anti-Western line of thought.

Sadly however I have not as yet found a Western, fully Roman Catholic, path to theosis... ... ..and I appreciate anyone who can point me in the proper direction.

\hfill

\texttt{Cologero on 2012-12-27 at 21:26 said:}

I believe, Jason-Adam, you are erecting artificial barriers. We pointed out that Scupoli's ``Spiritual Combat'' was adapted by the Orthodox with minor changes; we mentioned that for a reason. Furthermore, perhaps even others will gain from working through the exercises in that book, even if they need to make certain concessions in the language.

As for a Western path to theosis, there is plenty to choose from, but I'll list the Divine Comedy, the works of John of the Cross, ``The Three Ages of the Spiritual Life'' by Garrigou-Lagrange, and ``Mystical Evolution'' by Fr. Arintero.

catharsis : theoria : theosis = purgation : illumination : unification

As for hesychasm, I don't believe it is something anyone should be trying on his own.

\hfill

\texttt{Jason-Adam on 2012-12-27 at 22:25 said:}

I spoke too broadly and I ask for forgiveness for my error : Spiritual Combat/Unseen Warfare is indeed a classic of both Western and Eastern spirituality... ... What I was trying to say is that the theology of Gregory Palamas, which is the basis of Orthodox mysticism, is based on dividing the essence and energy of God, something the Church of Rome denies, thus despite whatever has been said in recent times one can not uphold Palamas and be Catholic in good faith. I am not denying the validity of Palamas for the Orthodox church, but I dont know if he can be accepted by Westerners as a guide.

\hfill

\texttt{Logres on 2012-12-27 at 22:45 said:}

Energy vs. Essence : I know it is a loaded discussion, but it does seem to preserve the idea of God's unmanifested Divinity, which (as I understand it so far) is central to Traditional thought in general.

\hfill

\texttt{Graham on 2012-12-28 at 14:35 said:}

``I believe, Jason-Adam, you are erecting artificial barriers.''

On this note, Jason-Adam, one might benefit from reading Marco Pallis' essay `The Catholic Church in Crisis'. I am certain it would help you orient yourself, and may indicate to you some areas in which the West could learn from the East.

``As for hesychasm, I don't believe it is something anyone should be trying on his own.''

Certain remarks of Guenon's on Yoga apply, especially since what has been published on hesychasm is, like Yoga, largely a popularization, as I indicated in an earlier comment:

``Nevertheless it must be added that these same `practices' can also have repercussions in the subtle modalities of the individual unsuspected by the ignorant person who undertakes them ... '' (Reign of Quantity, Confusion of the Psychic and the Spiritual)

\hfill

\texttt{Jason-Adam on 2012-12-30 at 18:44 said:}

I read Pallis' article but am not sure how far it can be accepted -- one of the errors condemned by B Pius IX was that of saying the Papacy was at fault in its dealing with the Eastern church... ... 

\hfill

\texttt{Logres on 2012-12-30 at 23:19 said:}

Graham,this is utterly right, I was thinking more of the lack of appreciation in the Church for the pagan foundations that go into Christianity, and which are an essential element of it, at least if it is to survive. The dialogues of later years have been aimed more at modernizing, and not even at neo-paganizing, certainly not at making explicit the sacerdotal and cultural foundations of the Faith. How to reconcile Boniface with St. Patrick, is more of what I was wondering?


\hfill

\end{sffamily}
\end{footnotesize}

